\subsection{各頂点は3分岐以下であること}

\begin{theorem}
  \label{thm:pos_degree_le_three}
  \lean{pos_degree_le_three}
  \leanok
  \uses{lem:symmMatrix_quadform_upper_bound}
  $C$はGCMであるとする．
  また，$S(C)$は正定値であるとする．
  このとき，任意の頂点$i$に対し，$i$の分岐数(生えている辺の数)は高々$3$である：
  \begin{equation}
    \deg(i) \le 3.
  \end{equation}
\end{theorem}

\begin{proof}
  $\# N(i) = \deg(i) \ge 4$と仮定して矛盾を導く．
  $x := \transpose{(x_1, x_2, \ldots, x_n)} \in \R^n$を以下で定める：
  \begin{equation}
    x_j = \begin{cases}
      2, & j = i,\\
      1, & j \in N(i),\\
      0, & \textrm{otherwise}.
    \end{cases}
  \end{equation}
  また，$E = \set{(i, a)}{a \in N(i)} \subseteq I \times I$と定める．

  補題~\ref{lem:symmMatrix_quadform_upper_bound}が使いたい．
  各条件を確認する．
  定義から明らかに，任意の$i$に対し，$0 \le x_i$が成り立つ．
  $(i, a) \in E$をとると，$E$の定め方から$a \in N(i)$である．
  また，$(a, i) \in E$と仮定すると，$i = a$となり$a \in N(i)$に矛盾するから$(a, i) \notin E$である．

  以上より補題~\ref{lem:symmMatrix_quadform_upper_bound}が使える．
  よって，
  \begin{align}
    \transpose{x} S(C) x
    &\le \sum_{i \in I} 2 x_i ^ 2 - 2 \sum_{p \in E} x_{p_1} x_{p_2} \\
    &= \sum_{i \in I} 2 x_i ^ 2 - 2 \cdot (2 \cdot \# N(i)) \\
    &= (8 + 2 \cdot \# N(i)) - 2 \cdot (2 \cdot \# N(i)) \\
    &= 8 - 2 \cdot \# N(i) \\
    &\le 0
  \end{align}
  となる．
  ただし，1つ目の等号は，$E$の定義と$x_i x_a = 2$より
  \begin{equation}
    \sum_{p \in E} x_{p_1} x_{p_2}
    = \sum_{a \in N(i)} x_{i} x_{a}
    = 2 \cdot (2 \cdot \# N(i))
  \end{equation}
  から従う．
  2つ目の等号は
  \begin{align}
    \sum_{i \in I} 2 x_i ^ 2
    &= \sum_{j = i} 2 x_j ^ 2 + \sum_{\substack{j \ne i \\ j \in N(i)}} 2 x_j ^ 2 \\
    &= \# \set{j}{j = i} \cdot (2 \cdot 2 ^ 2) + \# (\set{j}{i \ne j} \cap N(i)) \cdot (2 \cdot 1 ^ 2) \\
    &= 8 + 2 \cdot \# N(i)
  \end{align}
  から従う．

  一方，$S(C)$は正定値であり，$x \ne 0$であるから，
  \begin{equation}
    0 < \transpose{x} S(C) x
  \end{equation}
  である．
  よって，矛盾が生じたから$\deg(i) \le 3$が示せた．
\end{proof}